%!TeX root=../treasuretop.tex
\ArtChapter[The Ebb-Tide Runs]{chap23head}

   \lettrine[lines=4,image=true]{chap23initT}{he} coracle---as I had ample reason to know before I was done with her---was a very safe boat for a person of my height and weight, both buoyant and clever in a seaway; but she was the most cross-grained, lop-sided craft to manage. Do as you pleased, she always made more leeway than anything else, and turning round and round was the manoeuvre she was best at. Even Ben Gunn himself has admitted that she was \enquote{queer to handle till you knew her way.}

Certainly I did not know her way. She turned in every direction but the one I was bound to go; the most part of the time we were broadside on, and I am very sure I never should have made the ship at all but for the tide. By good fortune, paddle as I pleased, the tide was still sweeping me down; and there lay the \textit{Hispaniola} right in the fairway, hardly to be missed.

First she loomed before me like a blot of something yet blacker than darkness, then her spars and hull began to take shape, and the next moment, as it seemed (for, the farther I went, the brisker grew the current of the ebb), I was alongside of her hawser and had laid hold.

The hawser was as taut as a bowstring, and the current so strong she pulled upon her anchor. All round the hull, in the blackness, the rippling current bubbled and chattered like a little mountain stream. One cut with my sea-gully and the \textit{Hispaniola} would go humming down the tide.

So far so good, but it next occurred to my recollection that a taut hawser, suddenly cut, is a thing as dangerous as a kicking horse. Ten to one, if I were so foolhardy as to cut the \textit{Hispaniola} from her anchor, I and the coracle would be knocked clean out of the water.

This brought me to a full stop, and if fortune had not again particularly favoured me, I should have had to abandon my design. But the light airs which had begun blowing from the south-east and south had hauled round after nightfall into the south-west. Just while I was meditating, a puff came, caught the \textit{Hispaniola}, and forced her up into the current; and to my great joy, I felt the hawser slacken in my grasp, and the hand by which I held it dip for a second under water.

With that I made my mind up, took out my gully, opened it with my teeth, and cut one strand after another, till the vessel swung only by two. Then I lay quiet, waiting to sever these last when the strain should be once more lightened by a breath of wind.

All this time I had heard the sound of loud voices from the cabin, but to say truth, my mind had been so entirely taken up with other thoughts that I had scarcely given ear. Now, however, when I had nothing else to do, I began to pay more heed.

One I recognized for the coxswain’s, Israel Hands, that had been Flint’s gunner in former days. The other was, of course, my friend of the red night-cap. Both men were plainly the worse of drink, and they were still drinking, for even while I was listening, one of them, with a drunken cry, opened the stern window and threw out something, which I divined to be an empty bottle. But they were not only tipsy; it was plain that they were furiously angry. Oaths flew like hailstones, and every now and then there came forth such an explosion as I thought was sure to end in blows. But each time the quarrel passed off and the voices grumbled lower for a while, until the next crisis came and in its turn passed away without result.

On shore, I could see the glow of the great camp-fire burning warmly through the shore-side trees. Someone was singing, a dull, old, droning sailor’s song, with a droop and a quaver at the end of every verse, and seemingly no end to it at all but the patience of the singer. I had heard it on the voyage more than once and remembered these words:

\enquote{But one man of her crew alive,
What put to sea with seventy-five.}

And I thought it was a ditty rather too dolefully appropriate for a company that had met such cruel losses in the morning. But, indeed, from what I saw, all these buccaneers were as callous as the sea they sailed on.

At last the breeze came; the schooner sidled and drew nearer in the dark; I felt the hawser slacken once more, and with a good, tough effort, cut the last fibres through.

The breeze had but little action on the coracle, and I was almost instantly swept against the bows of the \textit{Hispaniola}. At the same time, the schooner began to turn upon her heel, spinning slowly, end for end, across the current.

I wrought like a fiend, for I expected every moment to be swamped; and since I found I could not push the coracle directly off, I now shoved straight astern. At length I was clear of my dangerous neighbour, and just as I gave the last impulsion, my hands came across a light cord that was trailing overboard across the stern bulwarks. Instantly I grasped it.

Why I should have done so I can hardly say. It was at first mere instinct, but once I had it in my hands and found it fast, curiosity began to get the upper hand, and I determined I should have one look through the cabin window.

I pulled in hand over hand on the cord, and when I judged myself near enough, rose at infinite risk to about half my height and thus commanded the roof and a slice of the interior of the cabin.

By this time the schooner and her little consort were gliding pretty swiftly through the water; indeed, we had already fetched up level with the camp-fire. The ship was talking, as sailors say, loudly, treading the innumerable ripples with an incessant weltering splash; and until I got my eye above the window-sill I could not comprehend why the watchmen had taken no alarm. One glance, however, was sufficient; and it was only one glance that I durst take from that unsteady skiff. It showed me Hands and his companion locked together in deadly wrestle, each with a hand upon the other’s throat.

   \begin{figure}[p]
\centering
\includegraphics[width=\linewidth]{chap23handswrestle}
\caption[Locked together in deadly wrestle]{Hands and his companion locked together in deadly wrestle}
\end{figure}  

I dropped upon the thwart again, none too soon, for I was near overboard. I could see nothing for the moment but these two furious, encrimsoned faces swaying together under the smoky lamp, and I shut my eyes to let them grow once more familiar with the darkness.

The endless ballad had come to an end at last, and the whole diminished company about the camp-fire had broken into the chorus I had heard so often:

\longyoho

I was just thinking how busy drink and the devil were at that very moment in the cabin of the \textit{Hispaniola}, when I was surprised by a sudden lurch of the coracle. At the same moment, she yawed sharply and seemed to change her course. The speed in the meantime had strangely increased.

I opened my eyes at once. All round me were little ripples, combing over with a sharp, bristling sound and slightly phosphorescent. The \textit{Hispaniola} herself, a few yards in whose wake I was still being whirled along, seemed to stagger in her course, and I saw her spars toss a little against the blackness of the night; nay, as I looked longer, I made sure she also was wheeling to the southward.

I glanced over my shoulder, and my heart jumped against my ribs. There, right behind me, was the glow of the camp-fire. The current had turned at right angles, sweeping round along with it the tall schooner and the little dancing coracle; ever quickening, ever bubbling higher, ever muttering louder, it went spinning through the narrows for the open sea.

Suddenly the schooner in front of me gave a violent yaw, turning, perhaps, through twenty degrees; and almost at the same moment one shout followed another from on board; I could hear feet pounding on the companion ladder and I knew that the two drunkards had at last been interrupted in their quarrel and awakened to a sense of their disaster.

I lay down flat in the bottom of that wretched skiff and devoutly recommended my spirit to its Maker. At the end of the straits, I made sure we must fall into some bar of raging breakers, where all my troubles would be ended speedily; and though I could, perhaps, bear to die, I could not bear to look upon my fate as it approached.

So I must have lain for hours, continually beaten to and fro upon the billows, now and again wetted with flying sprays, and never ceasing to expect death at the next plunge. Gradually weariness grew upon me; a numbness, an occasional stupor, fell upon my mind even in the midst of my terrors, until sleep at last supervened and in my sea-tossed coracle I lay and dreamed of home and the old \enquote{Admiral Benbow}.